% GENERATED FILE. Edit canonical lessons, then run npm run generate:books.
% canonical-source-hash: fnv1a64:8fbd5f718a5900f3
% canonical-lessons: TE-C196-phiju, TE-C196-jarimana, TE-C196-gaduvu, TE-C196-badhyata, TE-C196-darakhastu

\chapter{Fee, Fine, Deadline, Responsibility, Application}
\label{ch:te-fee-fine-deadline-responsibility-application}
\hlchaptermodality{196}

\begin{chapteropening}
\textbf{By the end of this chapter:} \emph{I can say a fee, a fine, a deadline, a responsibility and an application.}

Complete the last lesson of chapter 196: I can say a fee, a fine, a deadline, a responsibility and an application.
\end{chapteropening}

\section[\texorpdfstring{\te{ఫీజు} (phīju)}{ఫీజు (phīju)}]{\te{ఫీజు} (phīju) — a fee}

\label{lesson:TE-C196-phiju}



\begin{quote}
Before the new one: say the Telugu for a colleague, then the Telugu for a meeting.
\end{quote}

\subsection*{You'll want to know: \te{ఫీజు}}
\textbf{\te{ఫీజు}} — \emph{phīju} — \textquotedblleft{}a fee\textquotedblright{}.

\begin{grammarlens}[title={What you've built}]
One more word for this chapter.
\end{grammarlens}

\subsection*{Your turn}
\begin{itemize}
\raggedright
  \item \emph{Say it:} \emph{phīju}
  \item \emph{Say it:} \emph{phīju}, once more
  \item \emph{Say it:} say \emph{phīju}
  \item \emph{Recall:} say the Telugu for a colleague, then the Telugu for a meeting, then say \emph{phīju} again
\end{itemize}

\begin{tcolorbox}[breakable,colback=teal!4,colframe=teal!35!black,title={Before you move on}]
What does \te{ఫీజు} mean? (\textquotedblleft{}A fee\textquotedblright{}.) Say it once more.
\end{tcolorbox}
\section[\texorpdfstring{\te{జరిమానా} (jarimānā)}{జరిమానా (jarimānā)}]{\te{జరిమానా} (jarimānā) — a fine, a penalty}

\label{lesson:TE-C196-jarimana}



\begin{quote}
Before the new one: say the Telugu for a meeting, then the Telugu for a fee.
\end{quote}

\subsection*{You'll want to know: \te{జరిమానా}}
\textbf{\te{జరిమానా}} — \emph{jarimānā} — \textquotedblleft{}a fine, a penalty\textquotedblright{}.

\begin{grammarlens}[title={What you've built}]
One more word for this chapter.
\end{grammarlens}

\subsection*{Your turn}
\begin{itemize}
\raggedright
  \item \emph{Say it:} \emph{jarimānā}
  \item \emph{Say it:} \emph{jarimānā}, once more
  \item \emph{Say it:} say \emph{jarimānā}
  \item \emph{Recall:} say the Telugu for a meeting, then the Telugu for a fee, then say \emph{jarimānā} again
\end{itemize}

\begin{tcolorbox}[breakable,colback=teal!4,colframe=teal!35!black,title={Before you move on}]
What does \te{జరిమానా} mean? (\textquotedblleft{}A fine, a penalty\textquotedblright{}.) Say it once more.
\end{tcolorbox}
\section[\texorpdfstring{\te{గడువు} (gaḍuvu)}{గడువు (gaḍuvu)}]{\te{గడువు} (gaḍuvu) — a deadline, a time limit}

\label{lesson:TE-C196-gaduvu}



\begin{quote}
Before the new one: say the Telugu for a fee, then the Telugu for a fine.
\end{quote}

\subsection*{You'll want to know: \te{గడువు}}
\textbf{\te{గడువు}} — \emph{gaḍuvu} — \textquotedblleft{}a deadline, a time limit\textquotedblright{}.

\begin{grammarlens}[title={What you've built}]
One more word for this chapter.
\end{grammarlens}

\subsection*{Your turn}
\begin{itemize}
\raggedright
  \item \emph{Say it:} \emph{gaḍuvu}
  \item \emph{Say it:} \emph{gaḍuvu}, once more
  \item \emph{Say it:} say \emph{gaḍuvu}
  \item \emph{Recall:} say the Telugu for a fee, then the Telugu for a fine, then say \emph{gaḍuvu} again
\end{itemize}

\begin{tcolorbox}[breakable,colback=teal!4,colframe=teal!35!black,title={Before you move on}]
What does \te{గడువు} mean? (\textquotedblleft{}A deadline, a time limit\textquotedblright{}.) Say it once more.
\end{tcolorbox}
\section[\texorpdfstring{\te{బాధ్యత} (bādhyata)}{బాధ్యత (bādhyata)}]{\te{బాధ్యత} (bādhyata) — a responsibility, a duty}

\label{lesson:TE-C196-badhyata}



\begin{quote}
Before the new one: say the Telugu for a fine, then the Telugu for a deadline.
\end{quote}

\subsection*{You'll want to know: \te{బాధ్యత}}
\textbf{\te{బాధ్యత}} — \emph{bādhyata} — \textquotedblleft{}a responsibility, a duty\textquotedblright{}.

\begin{grammarlens}[title={What you've built}]
One more word for this chapter.
\end{grammarlens}

\subsection*{Your turn}
\begin{itemize}
\raggedright
  \item \emph{Say it:} \emph{bādhyata}
  \item \emph{Say it:} \emph{bādhyata}, once more
  \item \emph{Say it:} say \emph{bādhyata}
  \item \emph{Recall:} say the Telugu for a fine, then the Telugu for a deadline, then say \emph{bādhyata} again
\end{itemize}

\begin{tcolorbox}[breakable,colback=teal!4,colframe=teal!35!black,title={Before you move on}]
What does \te{బాధ్యత} mean? (\textquotedblleft{}A responsibility, a duty\textquotedblright{}.) Say it once more.
\end{tcolorbox}
\section[\texorpdfstring{\te{దరఖాస్తు} (darakhāstu)}{దరఖాస్తు (darakhāstu)}]{\te{దరఖాస్తు} (darakhāstu) — an application}

\label{lesson:TE-C196-darakhastu}



\begin{quote}
Before the new one: say the Telugu for a deadline, then the Telugu for a responsibility.
\end{quote}

\subsection*{You'll want to know: \te{దరఖాస్తు}}
\textbf{\te{దరఖాస్తు}} — \emph{darakhāstu} — \textquotedblleft{}an application\textquotedblright{}.

\begin{grammarlens}[title={What you've built}]
One more word for this chapter.
\end{grammarlens}

\subsection*{Your turn}
\begin{itemize}
\raggedright
  \item \emph{Say it:} \emph{darakhāstu}
  \item \emph{Say it:} \emph{darakhāstu}, once more
  \item \emph{Say it:} say \emph{darakhāstu}
  \item \emph{Recall:} say the Telugu for a deadline, then the Telugu for a responsibility, then say \emph{darakhāstu} again
\end{itemize}

\begin{tcolorbox}[breakable,colback=teal!4,colframe=teal!35!black,title={Before you move on}]
What does \te{దరఖాస్తు} mean? (\textquotedblleft{}An application\textquotedblright{}.) Say it once more.
\end{tcolorbox}
